% Branch A: "expected" outcome. Select when the rule in README (A) holds.
% Preamble file: defines result sentences used in abstract, intro and discussion.
\newcommand{\AbstractResults}{Under need-dependent silence, response-pooled adaptation under-served low-response residents (visit representation \vRhoSix). Population pooling with record-based screening raised their representation to \vRhoSevenR\ and their welfare by \vDWs\ (95\% CI \vDWsCI) at equal service use, and exceeded the best static policy in population welfare, most when capacity was scarce. Uncapped pooling amplified interpretation errors, which weight caps mitigated, and explicit control reduced policy variability relative to an informed language-model planner.}
\newcommand{\IntroEvidence}{In a data-grounded simulation evaluated on hidden utilities, response-pooled adaptation allocated the low-response cohort a visit share well below its share of need. Population pooling and record-based screening largely corrected this at equal service use and improved population welfare over the best static policy, with the largest gains under scarce capacity. Weight caps limited the amplification of interpretation errors, and the explicit controller was markedly more stable than an informed language-model planner.}
\newcommand{\DiscussionResults}{Selective participation, not interpretation alone, determined who was served: the distortion vanished without need-dependent silence and grew with it. Record-based screening contributed most of the correction, and population pooling added an approximately independent share. These gains came without additional service volume and were largest when capacity was scarce. Correction was incomplete at the strongest silence and did not extend to dissatisfaction that leaves no trace in group composition or records.}
